%% parts/p05/reverse-mathematics.tex --- Part 5
\chapter{Reverse Mathematics}\label{sec:p05-reverse-mathematics}

\noindent\textit{Reader note.}  This separately authored survey introduces the
standard second-order calibration used in the independence catalog.  It does
not replace the earlier full/Henkin semantics chapter or the inherited
first-order arithmetic theories.  The program and subsystem equivalences are
sourced to \cite{Friedman1975,Simpson2009}.

Reverse mathematics asks which set-existence axioms are required to prove a
theorem already formalized over a weak base. The reversal direction is as
important as the forward proof: an implication from a subsystem to a theorem
does not by itself calibrate strength.

\begin{hbdefinition}{def:p05-rca0}{The base theory \(\mathsf{RCA}_0\)}
The two-sorted language has number variables, set-of-number variables, the
number symbols \(0,S,+,\times\), and membership \(n\in X\).  Arithmetical
formulas have no set quantifiers but may have number and set parameters.
\(\mathsf{RCA}_0\) consists of:
\begin{enumerate}
  \item the number-sort copy of the seven axioms of \(\mathsf Q\) in
        Table~\ref{tab:p05-q-pa-axioms}, together with set extensionality
        \[
        \forall X\,\forall Y
        \bigl(\forall n(n\in X\leftrightarrow n\in Y)\to X=Y\bigr);
        \]
  \item for every \(\Sigma^0_1\) formula \(\varphi(n,\bar a,\bar A)\), the
        induction axiom
        \[
        \bigl(\varphi(0,\bar a,\bar A)\land
        \forall n(\varphi(n,\bar a,\bar A)\to
                    \varphi(Sn,\bar a,\bar A))\bigr)
        \to\forall n\,\varphi(n,\bar a,\bar A);
        \]
  \item for each \(\Sigma^0_1\) formula
        \(\varphi(n,\bar a,\bar A)\) and \(\Pi^0_1\) formula
        \(\psi(n,\bar a,\bar A)\), the \(\Delta^0_1\)-comprehension axiom
        \[
        \forall n(\varphi(n,\bar a,\bar A)\leftrightarrow
                   \psi(n,\bar a,\bar A))
        \to
        \exists X\,\forall n
          (n\in X\leftrightarrow\varphi(n,\bar a,\bar A)),
        \]
        where the comprehended \(X\) is not free in \(\varphi\) or \(\psi\).
\end{enumerate}
All displayed schemes permit the shown number and set parameters.  The name
is needed to fix this exact weak base rather than an informal notion of
computable mathematics.  It stands as a formal two-sorted theory, differs from
full second-order arithmetic by restricted induction and comprehension, and
is used as the common base for the reversals below \cite{Simpson2009}.
\index{RCA0@\(\mathsf{RCA}_0\)}
\end{hbdefinition}

\begin{hbdefinition}{def:p05-z2}{Full second-order arithmetic \(\mathsf Z_2\)}
In the language of Definition~\ref{def:p05-rca0}, \(\mathsf Z_2\) is
\(\mathsf{RCA}_0\) augmented by the full induction scheme for every
two-sorted formula and the full comprehension scheme
\[
 \exists X\,\forall n\,(n\in X\leftrightarrow\varphi(n))
\]
for every formula \(\varphi\) in which \(X\) is not free.  The name is needed
as the common ambient theory whose proper subsystems are compared below.  It
stands as a formal theory, differs from \(\mathsf{RCA}_0\) by unrestricted
induction and comprehension, and is used only to state that ambient relation.
\end{hbdefinition}

\begin{hbdefinition}{def:p05-big-five}{The five standard subsystems}
Over \(\mathsf{RCA}_0\), define:
\begin{itemize}
  \item \(\mathsf{WKL}_0\) adds weak K\"onig's lemma: for every coded tree
        \(T\subseteq2^{<\N}\), if
        \(\forall n\,\exists s\in T\,(|s|=n)\), then there is
        \(P\subseteq\N\) such that \(P\upharpoonright n\in T\) for every
        \(n\);
  \item \(\mathsf{ACA}_0\) is the theory fixed in
        Definition~\ref{def:p05-full-aca}; equivalently, over
        \(\mathsf{RCA}_0\) it adds comprehension for every arithmetical
        formula, with number and set parameters;
  \item \(\mathsf{ATR}_0\) adds arithmetical transfinite recursion: if
        \((D,\prec)\) is a coded countable well-order and \(\theta\) is
        arithmetical, there is \(Y\subseteq D\times\N\) whose sections satisfy
        \[
          Y_i=\{n:\theta(i,n,Y^{<i},\bar A)\},
          \qquad
          Y^{<i}=\{\langle j,n\rangle:j\prec i
                         \text{ and }n\in Y_j\};
        \]
  \item \(\Pi^1_1\text{-}\mathsf{CA}_0\) adds the same comprehension axiom
        for every \(\Pi^1_1\) formula, allowing universal set quantifiers over
        an arithmetical matrix.
\end{itemize}
The collective label “Big Five” is needed only as a retrieval name for this
specific list, including the base \(\mathsf{RCA}_0\); it stands for five
distinct theories, not a sixth theory.  The added principles are
nonredundant, each theory is a subsystem of \(\mathsf Z_2\), and the list is
used in the calibration table and inclusion figure \cite{Simpson2009}.
\index{WKL0@\(\mathsf{WKL}_0\)}\index{ACA0@\(\mathsf{ACA}_0\)}
\index{ATR0@\(\mathsf{ATR}_0\)}\index{Pi11CA0@\(\Pi^1_1\text{-}\mathsf{CA}_0\)}
\end{hbdefinition}

\begingroup\small
\begin{longtable}{@{}p{0.15\textwidth}p{0.28\textwidth}p{0.39\textwidth}p{0.10\textwidth}@{}}
\caption{Reverse-mathematics calibration.  Each nonbase row
states one exact equivalence over \(\mathsf{RCA}_0\); the first row identifies
the common base and deliberately makes no self-reversal claim.}
\label{tab:p05-reverse-calibration}\\
\toprule
Subsystem & Defining added scheme or principle & Representative theorem
equivalent over \(\mathsf{RCA}_0\) & Source \\
\midrule
\endfirsthead
\caption[]{Reverse-mathematics calibration.  Each nonbase row
states one exact equivalence over \(\mathsf{RCA}_0\); the first row identifies
the common base and deliberately makes no self-reversal claim (continued)}\\
\toprule
Subsystem & Defining added scheme or principle & Representative theorem
equivalent over \(\mathsf{RCA}_0\) & Source \\
\midrule
\endhead
\midrule\multicolumn{4}{r}{\emph{continued on next page}}\\
\endfoot
\endlastfoot
\(\mathsf{RCA}_0\) &
\(\Sigma^0_1\)-induction and \(\Delta^0_1\)-comprehension as displayed in
Definition~\ref{def:p05-rca0}. &
Common base for the other four reversals; no equivalence over itself is
asserted. & \cite{Simpson2009} \\
\addlinespace
\(\mathsf{WKL}_0\) &
Every infinite coded subtree of \(2^{<\N}\) has a path. &
Every countable cover of \([0,1]\) by coded rational open intervals has a
finite subcover. & \cite{Simpson2009} \\
\addlinespace
\(\mathsf{ACA}_0\) &
Comprehension for every arithmetical formula, with number and set parameters. &
For every injection \(f:\N\to\N\), the range
\(\{f(n):n\in\N\}\) exists as a set. & \cite{Simpson2009} \\
\addlinespace
\(\mathsf{ATR}_0\) &
Arithmetical transfinite recursion along every coded countable well-order. &
For any two coded countable well-orders, one is isomorphic to an initial
segment of the other. & \cite{Simpson2009} \\
\addlinespace
\(\Pi^1_1\text{-}\mathsf{CA}_0\) &
Comprehension for every \(\Pi^1_1\) formula. &
For every coded sequence \(\langle T_n:n\in\N\rangle\) of trees on
\(\N^{<\N}\), the set \(\{n:T_n\text{ is well-founded}\}\) exists. &
\cite{Simpson2009} \\
\bottomrule
\end{longtable}
\endgroup

Each equivalence includes the displayed coding of its mathematical objects in
second-order arithmetic.  For the Heine--Borel row, “countable cover” means a
sequence of rational intervals whose union contains every coded real in
\([0,1]\); for the well-order row, “initial segment” includes the whole target
as a possible segment.  Changing from countable coded objects to arbitrary
uncountable spaces changes the statement and may change its strength.

\begin{hbtheorem}{thm:p05-big-five-strict}{cited=Simpson2009}{Strictness of the standard ladder}
Over the usual metatheory, the implications
\[
 \mathsf{RCA}_0
 <\mathsf{WKL}_0
 <\mathsf{ACA}_0
 <\mathsf{ATR}_0
 <\Pi^1_1\text{-}\mathsf{CA}_0
\]
are strict with respect to theorem inclusion in the common language. The
strictness proofs use model constructions and conservation/ordinal methods;
the display is not a claim that every ordinary theorem falls at one of these
five levels.
\end{hbtheorem}

\begin{figure}[htbp]
\centering
\begin{tikzpicture}[
  theory/.style={draw,rounded corners,align=center,text width=.72\textwidth,
                 inner sep=5pt,font=\small},
  strict/.style={-{Stealth[length=2.2mm]},thick}
]
\node[theory] (rca) at (0,0)
  {\(\mathsf{RCA}_0\)\\base: effective coding and computable comprehension};
\node[theory] (wkl) at (0,1.65)
  {\(\mathsf{WKL}_0\)\\coded Heine--Borel compactness};
\node[theory] (aca) at (0,3.30)
  {\(\mathsf{ACA}_0\)\\the range of every injection exists};
\node[theory] (atr) at (0,4.95)
  {\(\mathsf{ATR}_0\)\\comparability of coded countable well-orders};
\node[theory] (pica) at (0,6.60)
  {\(\Pi^1_1\text{-}\mathsf{CA}_0\)\\well-founded-tree index sets exist};
\draw[strict] (rca) -- node[right,font=\scriptsize] {strict implication} (wkl);
\draw[strict] (wkl) -- node[right,font=\scriptsize] {strict implication} (aca);
\draw[strict] (aca) -- node[right,font=\scriptsize] {strict implication} (atr);
\draw[strict] (atr) -- node[right,font=\scriptsize] {strict implication} (pica);
\end{tikzpicture}
\caption{The reverse-mathematics Big Five.  Each arrow means
that the upper theory proves every theorem of the lower theory and strictly
more in the common language.  The attached theorem is the exact representative
from Table~\ref{tab:p05-reverse-calibration}.  The diagram asserts neither that
an arrow is a cover in the lattice of theories nor that the five systems
exhaust natural reverse-mathematical strengths.}
\label{fig:p05-big-five}
\end{figure}

\begin{hbmethod}{def:p05-reversal-method}{A reversal proof}
To calibrate a theorem \(P\), first formalize its objects and hypotheses in the
language of second-order arithmetic. Prove \(T\vdash P\) for a candidate
subsystem \(T\). Then, over the base \(\mathsf{RCA}_0\), encode an arbitrary
instance of the characteristic principle of \(T\) into an instance of \(P\)
and recover the desired set from a solution. The method name is necessary to
distinguish a two-direction equivalence from an upper-bound proof. It stands as
a metamathematical proof pattern, differs through its recovery direction, and
is used to read every row of the table.
\index{reverse mathematics!reversal method}
\end{hbmethod}

The program has important exceptions and refinements. Principles may lie
strictly between the standard systems, form incomparable families, or require
nonstandard induction/comprehension combinations. Conservation theorems can
show that a stronger-looking extension proves no new sentences of a specified
class. Accordingly, the Big Five are a successful recurring classification,
not an axiom that exhausts mathematical strength.

\section{Further reading for Part V}
\label{sec:p05-further-reading}

\subsection*{Historical and primary sources}

For incompleteness and arithmetized provability, see
\cite{Godel1931,Rosser1936,Lob1955,Solovay1976,GuaspariSolovay1979}.
The model-theoretic landmarks are \cite{Skolem1934,Tennenbaum1959}, and the
decidable reducts begin with \cite{Presburger1929,Skolem1930Arithmetic,
FefermanVaught1959}.  Natural arithmetic independence is developed in
\cite{ParisHarrington1977,KanamoriMcAloon1987,Goodstein1944,KirbyParis1982}.
Friedman's programmatic source for reverse mathematics is
\cite{Friedman1975}.

\subsection*{Modern systematic treatments}

For arithmetic coding, incompleteness, and models of arithmetic, see
\cite{Shoenfield1967,BoolosBurgessJeffrey2007,Kaye1991}.  Provability logic is
treated in \cite{Boolos1993}; reverse mathematics in \cite{Simpson2009}.  These
modern sources supply contemporary formulations without replacing the
historical attribution of the preceding results.
